\documentclass[11pt, a4paper, openany]{scrartcl}
\usepackage{xcolor}
\usepackage[colorlinks=true,allcolors=blue,pdfborder={0 0 0},pdfstartview={FitV},breaklinks,linktocpage]{hyperref}
\usepackage{color,amsmath,amsfonts,amssymb}
\usepackage[utf8]{inputenc}
\usepackage{enumitem}%für mehrspaltige Aufzählungen
\usepackage[ngerman, english]{babel}
\usepackage{epsfig}
\usepackage{pstricks}
\usepackage{graphics}
\graphicspath{{Figures/}}
\usepackage{bbm}
%\usepackage{showlabels}
\usepackage[labelfont=bf]{caption}
%\usepackage{nonfloat}
\usepackage{booktabs}%für Tabellen
\usepackage{pgfplots} %für Plots
\pgfplotsset{
compat=1.5,
tick label style={font=\small},
every axis/.append style={line width=0.8pt, tick style={line width=0.7pt}}
} %für richtigen Abstand der labels
\usetikzlibrary{shapes,arrows,fit,calc,positioning,arrows}
\usepackage[
  textheight=25cm,
left=2.3cm,
right=2.3cm,
headheight=25pt,
  includehead,includefoot,
  heightrounded,
]{geometry}
\makeatletter
\def\input@path{{./Figures/}}
\usepackage{fancyhdr}
\pagestyle{fancy}
\fancyfoot{}
\fancyhead{}
% \renewcommand{\chaptermark}[1]
% {\markboth{{\chaptername\ \thechapter:\  #1}}{}}
% \renewcommand{\sectionmark}[1]
% {{\markright{ \thesection\ #1}}} 
%\renewcommand{\sectionmark}[1]{\markboth{#1}{}}

%\renewcommand\sectionmark[1]{\markright{#1}}
%\renewcommand{\thesection}{\arabic{section}.}
\fancyhead[L]{}
\fancyhead[R]{\thepage}
\fancyhead[C]{\nouppercase{\leftmark}}
%\fancyhead[C]{\nouppercase{\rightmark}}


%\usepackage{sectsty}
\usepackage{url}
%\usepackage{breakurl}
\usepackage{amsthm}
\usepackage{graphicx}
\usepackage{tikzpagenodes}

\definecolor{darkgreen}{rgb}{0.0,0.6,0.0}
\definecolor{shadedgreen}{rgb}{0.8,0.95,0.8}
\definecolor{shadedblue}{rgb}{0.85,0.85,1.0}
\newlength\figureheight 
\newlength\figurewidth


\usepackage[varg]{txfonts} %font/Schrift
\usepackage{csquotes}
\usepackage[minnames=3,maxnames=30,backend=biber,natbib=true,sorting=none,citestyle=numeric-comp,giveninits=true]{biblatex}
\addbibresource{Literature.bib}
%\setlength\bibitemsep{0.5em}
\setlength{\biblabelsep}{0.4em}

\DeclareFieldFormat{pages}{#1}
\DeclareFieldFormat{journaltitle}{#1}
\DeclareFieldFormat{eprint}{#1}
\DeclareFieldFormat[article]{title}{\it#1\isdot}

\DeclareBibliographyDriver{article}{
\printnames{author}, \printfield{title},
\href{http://dx.doi.org/\thefield{doi}}
{\printfield{journaltitle} 
\textbf{\printfield{volume}}, \printfield{pages} 
(\printfield{year})}.}

\DeclareBibliographyDriver{misc}{
\printnames{author}, \printfield{title},
\href{https://arxiv.org/abs/\thefield{eprint}}{
arXiv preprint, arXiv:\printfield{eprint}  (\printfield{year})}.}

\DeclareBibliographyDriver{inbook}{\hspace{-0.1cm}
\printnames{author},
\printfield{title},
in \printfield{booktitle}.
\printlist{publisher}
(\printfield{year}).
}


\renewbibmacro*{name:andothers}{% Based on name:andothers from biblatex.def
  \ifboolexpr{
    test {\ifnumequal{\value{listcount}}{\value{liststop}}}
    and
    test \ifmorenames
  }
    {\ifnumgreater{\value{liststop}}{1}
       {\finalandcomma}
       {}%
     \andothersdelim\bibstring[\emph]{andothers}}
    {}}

\setenumerate[1]{label=[\arabic*],ref=\arabic*}
\newcommand{\paper}[4]{\item #1, \,\textit{#2}, \,\href{#3}{#4}.\\[-1.4em]}


\newcommand{\bE}{\mathbf{E}}

\begin{document}

\pagenumbering{Roman}

\begin{titlepage}\pdfbookmark[0]{Title}{Title}
  \sffamily
  \begin{center}
{
\includegraphics[width=9cm]{logo.pdf}
\\[10em]
\Huge \bfseries Output summary of the CUED program}
\\[5em]\large
The CUED program is developed and maintained by:
\\[5em]
Computational Condensed Matter Theory Group
\\[0.5em]
Head: Ferdinand Evers
  \\[0.5em]
Institute of Theoretical Physics
  \\[0.5em]
University of Regensburg
  \\[0.5em]
Universitätsstraße 31
  \\[0.5em]
D\,-\,93053 Regensburg
  \\[0.5em]
Germany
\\[5em]
Version: Git commit GITCOMMITNUMBER, \today
\\[5em]
  \end{center}{\large
  With contributions from (in alphabetic order): 
  \\[1em]
  Jack Crewse
  \\[0.5em]
  Patrick Grössing
  \\[0.5em]
  Maximilian Nitsch
  \\[0.5em]
  Adrian Seith
  \\[0.5em]
  Jan Wilhelm
  }
\end{titlepage}







\pagenumbering{arabic}
\pagestyle{plain}




\pdfbookmark[0]{Contents}{Contents}
\tableofcontents


\pagestyle{fancy}

\section{Electric-field pulse}
The following electric driving field is employed in the simulation:
\begin{align}
    \bE(t)  = E(t)\,\hat{e}_\phi
    \,,\hspace{2em}
    E(t) = E_0\,    \sin\Big(2\pi f_0\,(1+f_\text{chirp} t) + \varphi\Big)\, e^{-t^2/(2\alpha)^2}\,.
    \label{e1}
\end{align}
With parameters:
\begin{itemize}
    \item PLACEHOLDER-EFIELD-DIRECTION
\end{itemize}
\begin{figure}[h!]
   \centering
%\setlength\figureheight{9cm} 
%\setlength\figurewidth{12cm}
\input{Efield.tikz}
\caption{Electric driving field .}
    \label{fig:Efield}
\end{figure}

\section{Brillouin zone and \textit{k}-point grid}
REPLACEMENTBZ


\section{Hamiltonian, bandstructure and dipoles}
REPLACEMENTHAMILTONIAN


\section{Time evolution of the density matrix}

\section{Time-dependent currents}

\section{Frequency-resolved emission spectrum}

\section{Summary of input parameters}

\section{Timings}


\section{References}
\begin{enumerate}[leftmargin=*]

\paper{J.~Wilhelm, P.~Grössing, A.~Seith, J.~Crewse, M.~Nitsch, L.~Weigl, C.~Schmid, and F.~Evers}{Semi\-con\-duc\-tor-Bloch Formalism: Derivation and Application to High-Harmonic Generation from Dirac Fermions}{TOBEFILLED}{ 
Phys.~Rev.~B~\,\textbf{x}, y (2021)}
\label{Wilhelm2021}

\end{enumerate}

\end{document}
